% 第 3 章：统一 Newton-Raphson 潮流核心
% 本章文件由 \input 引入，不含导言区。
\section{统一 Newton--Raphson 潮流核心}\label{sec:newton}

\subsection{功能定位与入口}

\compmeta{engine::NewtonSolver}{\srcpath{include/hacdcpf/engine/solver/native/nle/newton\_solver.hpp}}

统一 Newton 求解器是全平台潮流的\textbf{生产主路径}：AC（$\theta,V_m$）、
DC（$V_{dc}$）与换流器耦合在同一方程组、同一雅可比中联立求解
（\textbf{不是} AC/DC 主从交替迭代）。入口签名
（include/hacdcpf/power\_flow/newton\_solver.hpp；实现 \srcpath{src/power\_flow/newton\_solver.cpp:solve}）：
\begin{quote}\footnotesize
\texttt{PowerFlowResult solve(const SolverData\& data, const PowerFlowOptions\& opt,\\
\hspace*{2em}const InitialState* init = nullptr) const;}
\end{quote}
可变成员 \texttt{PatternCache cache\_} 缓存雅可比上下文、稀疏模式与线性求解器实例
（见第~\ref{sec:overview}~章缓存体系）。每个求解器实例另复用一个
\texttt{SolverWorkspace}（include/hacdcpf/power\_flow/newton\_solver.hpp，接口
\srcpath{include/hacdcpf/power\_flow/workspace.hpp:SolverWorkspace}）：\texttt{vm}/\texttt{va}/\texttt{vdc}
经 \srcpath{prepare\_state} 按平启动重置、\texttt{mismatch}/\texttt{dx} 经
\srcpath{prepare\_equations} 按方程数重置，求解体内均为工作区成员引用
（src/power\_flow/newton\_solver.cpp:solve），而非每次求解在栈上新建；递归守卫
\texttt{workspace\_in\_use\_} 使同伦延拓的内层求解改用局部工作区，避免覆写外层状态
（:solve）。

\subsection{状态变量与节点类型}
未知量排序（\srcpath{build\_jacobian\_context}，src/power\_flow/newton\_solver.cpp:build\_jacobian\_context）：
\[
x=\big[\Delta\theta_{\,n_p};\ \Delta V_{m,\,n_q};\ \Delta V_{dc,\,n_{dc}}\big]^{\mathsf T},
\qquad n_{\mathrm{var}}=n_p+n_q+n_{dc}
\]
行/列映射（:build\_jacobian\_context）：非 slack 母线 $k$ 的有功行=角度列；PQ 母线的无功行=
$V_m$ 列（偏移 $n_p$）；非 slack DC 母线的 $P_{dc}$ 行=$V_{dc}$ 列（偏移 $n_p+n_q$）。
节点角色：
\begin{itemize}
  \item \textbf{Slack}：首个 SLACK 为主参考；\textbf{额外 SLACK 同样整体剔除}
        （$\theta$、$V_m$ 全固定），每个孤立 AC 区域各自保留参考（:build\_jacobian\_context）；
  \item \textbf{PV}：仅 P 方程、$\theta$ 未知，$V_m$ 固定；
  \item \textbf{PQ}：P、Q 方程与 $\theta$、$V_m$ 皆未知；
  \item \textbf{DC}：\texttt{DC\_V} 母线为 DC 参考（钉住）；\texttt{DC\_ISOLATED}
        剔除出方程组（:build\_jacobian\_context）；其余为 P 型 DC 母线。DC 参考由
        \srcpath{plan\_dc\_island\_references} 逐 DC 孤岛规划（:plan\_dc\_island\_references，
        详见第~\ref{sec:hybrid}~章）。
\end{itemize}
两种精确建模扩展（Direction 2/3）：\texttt{VDC\_VAC} 换流器交流母线在增强方程模式下
强制改 PQ、Q 行替换为 $V_m=V_{set}$（:build\_jacobian\_context）；半光滑 Newton 模式下
PV 全转 PQ 并记录 NCP 限值（:build\_jacobian\_context，见第~\ref{sec:robust}~章）。

\subsection{失配方程}
约定 $\mathrm{mismatch}=F^{\mathrm{spec}}-F^{\mathrm{calc}}$，解
$J\,\Delta x=\mathrm{mismatch}$，$J=\partial F^{\mathrm{calc}}/\partial x$
（注释见 \srcpath{src/power\_flow/jacobian\_builder.cpp:evaluate\_residual\_impl}）。
AC 节点注入（\srcpath{src/power\_flow/ac\_kernel.cpp:evaluate\_ac\_kernel\_serial}，
$\theta_{ij}=\theta_i-\theta_j$）：
\begin{align*}
P_i^{\mathrm{calc}}&=\sum_{j}V_iV_j\big(G_{ij}\cos\theta_{ij}+B_{ij}\sin\theta_{ij}\big)\\
Q_i^{\mathrm{calc}}&=\sum_{j}V_iV_j\big(G_{ij}\sin\theta_{ij}-B_{ij}\cos\theta_{ij}\big)
\end{align*}
指定注入含 ZIP 负荷（见第~\ref{sec:assembly}~章）；VSC 交流注入加入 spec。
DC 节点：$\mathbf p_{dc}^{\mathrm{calc}}=\mathbf v_{dc}\odot(G_{dc}\mathbf v_{dc})$，
失配行 $=p_{dc}^{\mathrm{spec}}-p_{dc}^{\mathrm{calc}}$（src/power\_flow/jacobian\_builder.cpp:evaluate\_residual\_impl）。
装配入口 \srcpath{evaluate\_residual\_and\_jacobian}（:evaluate\_residual\_and\_jacobian）与
\srcpath{evaluate\_residual\_only}（:evaluate\_residual\_only）共享 \srcpath{evaluate\_residual\_impl}（:evaluate\_residual\_impl），
返回 $\lVert\mathrm{mismatch}\rVert_\infty$（:evaluate\_residual\_impl）。

\subsection{雅可比结构}
$J=\begin{bmatrix}H&N\\ M&L\end{bmatrix}$ 为极坐标非归一化（MATPOWER）形式，
增量为 $\Delta V$ 本身（非 $\Delta V/V$）：
$H=\partial P/\partial\theta$、$N=\partial P/\partial V$、$M=\partial Q/\partial\theta$、
$L=\partial Q/\partial V$。

非对角元（$i\neq j$，AC kernel 累加，src/power\_flow/ac\_kernel.cpp:evaluate\_ac\_kernel\_serial）：
\begin{align*}
H_{ij}&\mathrel{+}=V_iV_j\big(G_{ij}\sin\theta_{ij}-B_{ij}\cos\theta_{ij}\big) &
M_{ij}&\mathrel{+}=-V_iV_j\big(G_{ij}\cos\theta_{ij}+B_{ij}\sin\theta_{ij}\big)\\
N_{ij}&\mathrel{+}=V_i\big(G_{ij}\cos\theta_{ij}+B_{ij}\sin\theta_{ij}\big) &
L_{ij}&\mathrel{+}=V_i\big(G_{ij}\sin\theta_{ij}-B_{ij}\cos\theta_{ij}\big)
\end{align*}
对角元（src/power\_flow/jacobian\_builder.cpp:evaluate\_residual\_impl），含 ZIP 负荷导数，
$\underline V=$ \fld{min\_vm\_pu}（默认 $10^{-8}$）：
\begin{align*}
H_{ii}&=-Q_i^{\mathrm{calc}}-B_{ii}V_i^2 &
M_{ii}&=P_i^{\mathrm{calc}}-G_{ii}V_i^2\\
N_{ii}&=\frac{P_i^{\mathrm{calc}}}{\max(|V_i|,\underline V)}+G_{ii}V_i+\big(p_dw_{p1}+2V_ip_dw_{p2}\big)\\
L_{ii}&=\frac{Q_i^{\mathrm{calc}}}{\max(|V_i|,\underline V)}-B_{ii}V_i+\big(q_dw_{q1}+2V_iq_dw_{q2}\big)
\end{align*}
DC 块：对角 $p_{dc,i}^{\mathrm{linear}}+g_{ii}v_{dc,i}$、非对角 $g_{ij}v_{dc,i}$
（:evaluate\_residual\_impl）。换流器导数分三档注入（详见第~\ref{sec:hybrid}~章）：DC 自洽导数
始终注入（:evaluate\_residual\_impl）；AC$\leftrightarrow$DC 交叉耦合块仅在
\fld{enable\_coupled\_jacobian} 时注入（:evaluate\_residual\_impl）；AC 构网换流器的能量导管
无条件生效（:evaluate\_residual\_impl）。

\subsection{稀疏模式与求值内核}
\begin{itemize}
  \item \srcpath{build\_jacobian\_pattern}（\srcpath{src/power\_flow/jacobian\_pattern.cpp:build\_jacobian\_pattern}）：
        从 $Y_{bus}$ 稀疏结构生成 $4\times nnz(Y)$ 个 AC 块槽位 + DC 块 + 耦合槽；
        以 \fld{entry\_to\_nz} 哈希把 (row,col)$\to$nz 下标写入
        \texttt{ACEntry\{p\_va\_nz, p\_vm\_nz, q\_va\_nz, q\_vm\_nz, g, b\}}（:build\_jacobian\_pattern）；
        模式固定后每次迭代仅 \texttt{valuePtr()} 清零再累加
        （src/power\_flow/jacobian\_builder.cpp:evaluate\_residual\_impl），不再分配；
  \item \srcpath{evaluate\_ac\_kernel\_serial}/\srcpath{\_parallel}
        （src/power\_flow/ac\_kernel.cpp:evaluate\_ac\_kernel\_serial）：遍历 \fld{ac\_entries} 同时累加
        $P^{\mathrm{calc}}$/$Q^{\mathrm{calc}}$ 与（可选）雅可比值；并行版每线程
        独立累加器再归约，条目数 $<1024$ 或单线程回退串行（:effective\_ac\_threads，
        \srcpath{util::ThreadPool::global()}），线程数由 \fld{ac\_eval\_threads} 控制；
  \item 一致性校验工具 \srcpath{check\_jacobian}（\srcpath{src/power\_flow/jacobian\_check.cpp:check\_jacobian}）：
        中心差分 $h=10^{-6}$、容差 $10^{-5}$。
\end{itemize}

\subsection{线性求解}
抽象 \srcpath{SparseLinearSolver}\{analyze\_pattern, factorize, solve, backend\_name\}
（\srcpath{include/hacdcpf/engine/kernel/linear\_algebra/linear\_solver.hpp}，转发 MIPSolvers）。
默认后端优先级 \textbf{KLU $>$ UMFPACK $>$ MKL Pardiso $>$ SuperLU $>$ Eigen SparseLU}
（\srcpath{MIPSolvers/src/engine/kernel/linear\_algebra/linear\_solver.cpp:make\_default\_sparse\_solver}；
CMake 选项 \srcpath{HACDCPF\_USE\_SUITESPARSE}，缺失回退 Eigen SparseLU）。
稀疏模式仅在缓存失效时 \texttt{analyzePattern} 一次；每次 Newton 迭代重新数值分解
（\srcpath{solve\_linear}，src/power\_flow/newton\_solver.cpp:solve\_linear（lambda）；计数
\fld{factorization\_calls}/\fld{linear\_solve\_calls}）。坏条件数
（代理 $\kappa=(\max/\min$ 行范数$)\times(\max/\min$ 列范数$)$，:estimate\_condition\_proxy）
可触发 NK-GMRES 回退（:solve\_linear\_nk（lambda），默认关闭，见第~\ref{sec:robust}~章）。

\subsection{更新、步长截断与收敛}
\begin{itemize}
  \item \textbf{步长截断} \srcpath{clip\_step}（:clip\_step）：
        $|\Delta\theta|\le 1.5$~rad、$|\Delta V_m|\le 0.5$~pu、$|\Delta V_{dc}|\le 0.5$~pu；
        $V_m,V_{dc}$ 下限 \texttt{kMinVm}$=0.05$（:kMinVm，:NewtonSolver::solve）；
  \item \textbf{收敛判据}：缩放后 $\lVert F\rVert_\infty<\texttt{opt.tol}$；
        \fld{tol}=$10^{-8}$、\fld{max\_iter}=50（\texttt{Defaults::kPFTol}/\texttt{kPFMaxIter}，
        \srcpath{include/hacdcpf/model/defaults.hpp:Defaults}）；NaN/Inf 立即判失败（src/power\_flow/newton\_solver.cpp:solve）；
  \item \textbf{诚实收敛}：最终以末态残差重新判定 \fld{out.converged}（:solve）；
        DC 电压越限仅告警（:solve）；
  \item \textbf{缩放}（默认全开）：$D_F=\mathrm{diag}(1/\max(|F_i^{\mathrm{spec}}|,1))$、
        $D_x=\mathrm{diag}(\pi,V_i,V_{dc,i})$，解 $\hat J\Delta\hat x=\hat F$，
        $\hat J=D_FJD_x^{-1}$（\srcpath{src/power\_flow/nonlinear\_scaling.cpp:build\_nonlinear\_scaling}，
        钳制 $[10^{-6},10^6]$；详见第~\ref{sec:robust}~章）。
\end{itemize}

\subsection{PV$\leftrightarrow$PQ 外环}
MATPOWER 风格无功限值切换（src/power\_flow/newton\_solver.cpp:solve）：外环上限 30 次；
内层收敛后检查隐含无功 $Q_g^{\mathrm{implied}}=Q^{\mathrm{calc}}+Q^{\mathrm{load}}-Q^{\mathrm{conv}}$
越限（迟滞 \fld{pv\_q\_hysteresis\_pu}=0.01），越限则转 PQ 并把 \fld{qg} 钉在限值；
全部限内后尝试恢复 PV（\fld{pv\_recover\_vm\_tol\_pu}=0.01）。切换\textbf{只在外环}
发起，迭代内不做切换：内层 Newton 收敛后统一检查越限（:solve）；内层未收敛时
仅当残差 $\le 0.1$ 才允许按当前 $Q$ 估计切换重试（:solve，守卫避免大残差下
按不可靠 $Q$ 值在 PV/PQ 间来回振荡）。
换流器 PQ$\leftrightarrow$VDC\_Q 模式切换在残差 $<10^{-6}$ 后按迟滞评估
（高/低阈值 0.03/0.01 pu、持续 2 次迭代，:apply\_converter\_mode\_switching，:NewtonSolver::solve；
逻辑见第~\ref{sec:hybrid}~章）。

\subsection{失败回退链与全局化}
默认全局化策略为线搜索（\fld{globalization}=LineSearch），回退链依次为
（src/power\_flow/newton\_solver.cpp:solve）：
\begin{enumerate}
  \item 非单调 Armijo-GLL 线搜索（\srcpath{run\_line\_search}，
        src/power\_flow/newton\_solver.cpp:run\_line\_search（lambda））：回溯 $\alpha\leftarrow\beta\alpha$、
        上限 \fld{max\_line\_search\_steps}=10 次试探（:run\_line\_search（lambda）），接受条件
        $\phi(x+\alpha\Delta x)\le\varphi_k^{\max}+c_1\,\alpha\,\nabla\phi(x)^{\mathsf T}\Delta x$
        （\srcpath{include/hacdcpf/power\_flow/nonmonotone\_linesearch.hpp}）；窗口 $M=5$、$c_1=10^{-4}$、
        $\beta=0.5$（\fld{nonmonotone\_window}/\fld{armijo\_c}/\fld{line\_search\_beta}，
        include/hacdcpf/power\_flow/robust\_nonlinear\_options.hpp），merit 每轮压入窗口
        （src/power\_flow/newton\_solver.cpp:solve）；\fld{enable\_nonmonotone\_linesearch}=false
        时以 $M=1$ 构造、退化为单调 Armijo（:run\_line\_search（lambda））；
  \item 对角正则化 $J+\lambda I$（$\lambda_0=10^{-6}$、增长 $\times10$、最多 8 步），
        每个正则化方向再经同一线搜索判定接受（:solve）；
  \item 自动恢复调度（\fld{enable\_auto\_fallback\_scheduling} 默认\textbf{开}，
        include/hacdcpf/power\_flow/robust\_nonlinear\_options.hpp）：LM 恢复
        $(J^{\mathsf T}J+\lambda I)\Delta x=-J^{\mathsf T}F$，最多 4 次尝试、
        $\lambda$ 自 \fld{lm\_lambda0} 起 $\times10$ 增至 \fld{lm\_lambda\_max}
        （src/power\_flow/newton\_solver.cpp:solve）；仍失败则以 $J+I/\delta t$ 做 PTC 恢复并
        无条件接受（:solve）；全部失败记 \fld{inner\_failed}（:solve）。
\end{enumerate}
另两种全局化策略：信赖域 dogleg（半径 1.0$\to$10）与伪瞬态延拓
（$J+I/\delta t$，SER 时间步由 \fld{ptc\_dt0} 族驱动，详见
第~\ref{sec:robust:lmptc}~节），
由 \fld{opt.globalization} 三选一。完整算法伪代码见技术笔记
\srcpath{docs/archive/reference/technical\_notebook/sections/03\_power\_flow.tex} 的 Algorithm PF-1；
五阶段鲁棒管线的组件细节见本手册第~\ref{sec:robust}~章。

\subsection{结构闭合检查与结果}
首次装配后扫描空行/空列（\srcpath{scan\_empty\_rows\_cols}，:scan\_empty\_rows\_cols），
做 $n_{eq}=n_{var}$ 闭合校核，结果入 \fld{diagnostics.equation\_closure\_*}。
输出 \texttt{PowerFlowResult\{vm, va, vdc, converged, iterations, residual,
profiling, diagnostics, branch\_flows, vsc\_transfers, ...\}}
（\srcpath{include/hacdcpf/power\_flow/power\_flow\_result.hpp:PowerFlowResult}），字段全集见第~\ref{sec:options}~章。
迭代计数 \fld{out.iterations} = 外层 + 内层 + 1（src/power\_flow/newton\_solver.cpp:solve）。

\subsection{验证与校核}
\begin{itemize}
  \item \textbf{MATPOWER 回归}：\srcpath{tests/test\_matpower\_cases.cpp} 三档
        85 算例；Tier1 九算例逐案断言（收敛 + $|V|\in[0.5,1.5]$、$|\delta|<60^\circ$、
        残差 $<10^{-4}$），Tier2 批收敛率 $\ge 85\%$，Tier3 大系统至约 8.2 万节点；
        \srcpath{tests/test\_powerflow\_performance.cpp} 用独立残差复算交叉验证；
  \item \textbf{雅可比对差分}：\srcpath{tests/test\_unified\_pf\_upgrade.cpp}
        解析雅可比 vs 中心差分 max\_rel\_error $<10^{-3}$（差分 $\varepsilon=10^{-6}$）；
  \item \textbf{缩放回归}：\srcpath{tests/test\_scaling\_regression.cpp}
        （case9/14/30/57/118 默认缩放开收敛、关缩放回退仍收敛）；
  \item \textbf{基准性能}：技术笔记 03 节 84 算例基准表（Apple M4 单线程）：
        NR 对 8.2 万节点 1.75~s/17 次迭代，实际标度约 $O(N^{1.2})$；
  \item 已知边界：PV$\leftrightarrow$PQ 为离散切换外环（半光滑 NCP 为可选替代）；
        换流器不等式约束仅事后校核（见第~\ref{sec:hybrid}~章诚实口径）。
\end{itemize}
